\documentclass[12pt,reqno]{amsart}

% Version 6. Based on the user-supplied v5, with an independently readable
% introduction containing the principal conclusions. Exact repeated theorem
% bodies are stored below as label-free macros in this standalone source.
% Literature baseline: v4 (22 September 2026); targeted additions checked
% on 24 September 2026. Source locators refer to the versions in the bibliography.
% Standalone source. Compile with pdfLaTeX (two runs) or latexmk -pdf.
\usepackage[T1]{fontenc}
\usepackage{amsmath,amssymb,amsthm,mathtools}
\IfFileExists{stix2.sty}{\usepackage{stix2}}{%
  \IfFileExists{lmodern.sty}{\usepackage{lmodern}}{}%
  \PackageWarningNoLine{dhym-survey}{STIX2 unavailable; using installed fallback fonts}%
}
\IfFileExists{microtype.sty}{\usepackage{microtype}}{}
\usepackage{geometry}
\geometry{a4paper,textwidth=6.1in,top=1.2in,bottom=1.2in,
  headheight=14pt,headsep=0.2in,footskip=0.5in}
\usepackage{enumitem,needspace}
\setlist[itemize]{leftmargin=*,topsep=0.35\baselineskip,
  itemsep=0.15\baselineskip,parsep=0pt}
\setlist[enumerate]{leftmargin=*,label=\textup{(\roman*)},
  topsep=0.35\baselineskip,itemsep=0.15\baselineskip,parsep=0pt}
\usepackage{hyperref}
\hypersetup{colorlinks=true,linkcolor=blue,citecolor=blue,urlcolor=blue,
  linktocpage=true,bookmarksnumbered=true,bookmarksdepth=2,
  pdftitle={Positivity and Solvability of the Scalar Deformed Hermitian-Yang-Mills Equation},
  pdfauthor={},pdfsubject={Version 6: numerical criteria, Thomas-Yau, flows and stability}}
\linespread{1.1}
\emergencystretch=1.2em
\widowpenalty=10000
\clubpenalty=10000
\pagestyle{plain}
\setcounter{tocdepth}{1}
\numberwithin{equation}{section}

\theoremstyle{plain}
\newtheorem{theorem}{Theorem}[section]
\newtheorem{lemma}[theorem]{Lemma}
\newcommand{\recallheading}{Theorem}
\newtheorem*{recalledtheorem}{\recallheading}
\newenvironment{theoremrecall}[1]{%
 \renewcommand{\recallheading}{Theorem~\ref{#1} (recalled)}%
 \begin{recalledtheorem}%
}{\end{recalledtheorem}}
\theoremstyle{definition}
\newtheorem{definition}[theorem]{Definition}
\newtheorem{example}[theorem]{Example}
\newtheorem{problem}[theorem]{Problem}

\newcommand{\ddbar}{\sqrt{-1}\partial\overline{\partial}}
\newcommand{\cK}{\mathcal K}
\newcommand{\cF}{\mathcal F}
\newcommand{\cH}{\mathcal H}
\newcommand{\R}{\mathbb R}
\newcommand{\C}{\mathbb C}
\newcommand{\bbP}{\mathbb P}
\DeclareMathOperator{\NS}{NS}

\title[Positivity and solvability of the dHYM equation]{Positivity and Solvability of the\\Scalar Deformed Hermitian--Yang--Mills Equation}
\date{Version 6 --- 24 September 2026}


% Canonical statement bodies. No labels, wrappers or external input files.
\newcommand{\SurfaceCriterionStatement}{%
Let $X$ be a compact K\"ahler surface, $\Omega\in\cK_X$ and
$A\in H^{1,1}(X,\R)$ with $\langle A\smile\Omega,[X]\rangle>0$.
Set
\[
 \widehat\theta=\operatorname{Arg}\langle(\Omega+\sqrt{-1}A)^2,[X]\rangle
 \in(0,\pi),\qquad c=\cot\widehat\theta.
\]
The following are equivalent:
\begin{enumerate}
\item For some, equivalently every, K\"ahler form $\omega\in\Omega$,
there is a smooth closed real $(1,1)$-form $\eta\in A$ with
$\vartheta_\omega(\eta)=\widehat\theta$.
\item $A+c\Omega\in\cK_X$.
\item $\langle A+c\Omega,[C]\rangle>0$ for every irreducible curve
$C\subset X$.
\end{enumerate}
}

\newcommand{\SubsolutionStatement}{%
Let $(X,\omega)$ be a compact K\"ahler $n$-fold with $n\geq2$, and
let $\alpha$ be a smooth closed real $(1,1)$-form. Fix
$\widehat\theta\in((n-2)\pi/2,n\pi/2)$. Assume both of the following:
\begin{itemize}
\item $e^{-\sqrt{-1}\widehat\theta}\int_X(\omega+\sqrt{-1}\alpha)^n
\in\R_{>0}$.
\item There is a smooth $C$-subsolution
$\underline\varphi\in C^\infty(X,\R)$ for the target $\widehat\theta$.
\end{itemize}
Then $\vartheta_\omega(\alpha+\ddbar\varphi)=\widehat\theta$ has a
smooth solution, unique modulo constants.
}

\newcommand{\RayCriterionStatement}{%
Let $X$ be a compact K\"ahler $n$-fold with $n\geq2$, let
$A\in H^{1,1}(X,\R)$ and let $\omega\in\Omega$ be a K\"ahler form.
Fix $\theta\in(0,\pi)$ and assume
$\langle\cF_n(A,\Omega;\theta),[X]\rangle=0$.
Then there is a smooth closed real $(1,1)$-form $\eta\in A$ satisfying
$Q_\omega(\eta)=\theta$ if and only if, for every $t\geq0$ and every
irreducible analytic subvariety
$V\subseteq X$ of dimension $1\leq p\leq n$,
\[
 \big\langle\cF_p(A+t\Omega,\Omega;\theta),[V]\big\rangle
 \begin{cases}
 >0,&p<n,\\
 \geq0,&p=n.
 \end{cases}
\]
}

\newcommand{\ProjectiveCriterionStatement}{%
Let $X$ be smooth projective of dimension $n\geq2$, let
$A\in H^{1,1}(X,\R)$ and let $\omega\in\Omega$ be a K\"ahler form.
Fix $\theta\in(0,\pi)$.
Then there is a smooth closed real $(1,1)$-form $\eta\in A$ satisfying
$Q_\omega(\eta)=\theta$ if and only if both of the following hold:
\begin{itemize}
\item $\langle\cF_n(A,\Omega;\theta),[X]\rangle=0$.
\item $\langle\cF_p(A,\Omega;\theta),[V]\rangle>0$ for every irreducible
proper analytic subvariety $V\subset X$ with $1\leq p=\dim V<n$.
\end{itemize}
}

\newcommand{\PropernessStatement}{%
Let $(X,\omega)$ be a compact K\"ahler $n$-fold with $n\geq4$, and
let $\alpha$ be a smooth closed real $(1,1)$-form. Fix
$0<\theta<\pi/2$ and $\widehat\theta=n\pi/2-\theta$. Assume
\begin{itemize}
\item $e^{-\sqrt{-1}\widehat\theta}\int_X(\omega+\sqrt{-1}\alpha)^n
\in\R_{>0}$;
\item $0\in\cH$.
\end{itemize}
Then $Q_\omega(\alpha_\varphi)=\theta$ is smoothly solvable if and
only if, for every $\varepsilon\in(0,\theta)$, there exist
$\delta_\varepsilon,C_\varepsilon>0$ such that
\[
 \mathcal J(\varphi)\geq\delta_\varepsilon d_1(\varphi,0)-C_\varepsilon
 \quad\text{for all }\varphi\in\cH_\varepsilon
 \text{ with }\sup_X\varphi=0.
\]
}

\newcommand{\ToricMirrorStatement}{%
Let $X$ be a smooth projective toric weak Fano $n$-fold, $n\geq2$,
$\omega\in\Omega$ torus invariant, and $M^\vee$ ample. Put
$A=c_1(M^\vee)$ and choose a compatible
$\widehat\theta\in((n-2)\pi/2,n\pi/2)$, with
$\theta=n\pi/2-\widehat\theta$. Assume
\[
 \langle\cF_p(A,\Omega;\theta),[V]\rangle\ne0
 \quad\text{for every proper toric }V,
 \quad 1\leq p=\dim V<n.
\]
For a sufficiently large integer $k$, let $\mathcal L_k$ be the SYZ
section associated with $M^{\otimes k}$ in the mirror of $(X,k\omega)$.
There are test objects $\mathcal E_{V,k}$, indexed by proper toric
subvarieties $V$, and morphisms
\[
 \mathcal E_{V,k}[-r]\longrightarrow\mathcal L_k,
 \qquad r=\operatorname{codim}_X V,
\]
such that $\mathcal L_k$ has a supercritical special SYZ representative
in its Hamiltonian class if and only if
\[
 \operatorname{Im}\frac{(-1)^rP_k(\mathcal E_{V,k})}
                          {P_k(\mathcal L_k)}<0
 \qquad\text{for every such }V.
\]
}

\newcommand{\SurfaceStabilityStatement}{%
Let $X$ be a smooth projective surface, $\Omega\in\NS(X)_\R$ an ample
class, and $L$ a line bundle. Put $A=c_1(L)$ and assume
$s=\langle A\smile\Omega,[X]\rangle>0$. Set
\[
 A_c=A+\frac{\langle\Omega^2-A^2,[X]\rangle}{2s}\Omega.
\]
The following are equivalent:
\begin{enumerate}
\item $\langle A_c,[C]\rangle\geq0$ for every integral curve $C\subset X$.
\item $L^{\otimes k}$ is $\sigma_{0,k\Omega}$-stable for every integer $k\geq1$.
\item $L^{\otimes k}$ is $\sigma_{0,k\Omega}$-stable for all sufficiently
large integers $k$.
\end{enumerate}
}

\begin{document}
\begin{abstract}
In the positive supercritical range, smooth scalar deformed
Hermitian--Yang--Mills solvability is characterized by a shifted
K\"ahler class on surfaces, by subsolutions analytically, and by
intersection inequalities along a ray on compact K\"ahler manifolds.
For projective manifolds, only the endpoint inequalities are needed.
In the stated hypercritical range, properness is another equivalent
criterion. We give the distinct hypotheses for three scalar flows,
separate special Lagrangian existence from smooth-flow convergence in
the Thomas--Yau problem, and state a toric large-volume correspondence.
On surfaces, stability under every integral scaling allows a nef
boundary class and therefore does not imply smooth dHYM solvability.
\end{abstract}
\maketitle
\tableofcontents

\section{Introduction}\label{background-forms-and-existence}
\subsection{The equation and its branch}
Let $X$ be a compact K\"ahler manifold of dimension $n\geq2$.
Fix a real class $A\in H^{1,1}(X,\R)$, a K\"ahler class
$\Omega\in\cK_X$, and representatives $\alpha\in A$ and
$\omega\in\Omega$, with $\alpha$ smooth and closed and $\omega>0$.
For real smooth potentials, put
\begin{equation}\label{potentials-and-cohomology-classes}
 \alpha_\varphi=\alpha+\ddbar\varphi,\qquad
 \omega_\psi=\omega+\ddbar\psi>0.
\end{equation}
For a K\"ahler form $\chi$ and a real $(1,1)$-form $\eta$, let
$\lambda_j(\eta;\chi)$ be the eigenvalues of $\eta$ relative to $\chi$.
The phase and its determinant identity are
\begin{align}
 \vartheta_\chi(\eta)
 &=\sum_{j=1}^n\arctan\lambda_j(\eta;\chi),
 \qquad \arctan\lambda_j\in(-\pi/2,\pi/2),
 \label{scalar-phase-operator}\\
 (\chi+\sqrt{-1}\eta)^n
 &=e^{\sqrt{-1}\vartheta_\chi(\eta)}
   \prod_{j=1}^n\sqrt{1+\lambda_j(\eta;\chi)^2}\,\chi^n.
 \label{eigenvalue-identity}
\end{align}
Fix a \emph{supercritical} target
$\widehat\theta\in((n-2)\pi/2,n\pi/2)$.
The scalar deformed Hermitian--Yang--Mills equation asks that this
real phase be constant:
\begin{itemize}
\item For fixed $\omega$, the unknown is
$\varphi\in C^\infty(X,\R)$, and
\begin{equation}\label{dHYM-equation}
 \vartheta_\omega(\alpha_\varphi)=\widehat\theta.
\end{equation}
\item With only $(A,\Omega)$ fixed, the unknown is
$(\varphi,\psi)\in C^\infty(X,\R)^2$, and
\begin{equation}\label{dHYM-with-varying-background}
 \vartheta_{\omega_\psi}(\alpha_\varphi)=\widehat\theta,
 \qquad \omega_\psi>0.
\end{equation}
\end{itemize}
Set
\begin{equation}\label{central-charge}
 Z_X(A,\Omega)=\big\langle(\Omega+\sqrt{-1}A)^n,[X]\big\rangle.
\end{equation}
A solution satisfies
\begin{equation}\label{phase-compatibility}
 e^{-\sqrt{-1}\widehat\theta}Z_X(A,\Omega)
 =\int_X\prod_{j=1}^n
   \sqrt{1+\lambda_j(\alpha_\varphi;\omega_\psi)^2}\,
   \omega_\psi^n>0,
\end{equation}
with $\psi=0$ for a fixed background. We call this the compatibility
condition. It is imposed wherever a compatible target is specified.

\subsection{Existence: surfaces, subsolutions and numerical tests}
\label{intro-existence-results}
\begin{theorem}[Jacob--Yau {\cite[Theorem 1.2]{JY17}}]
\label{thm-surface-criterion}
\SurfaceCriterionStatement
\end{theorem}
The proof in Section~\ref{surface-criterion} reduces the equation to
\[
 \beta_\varphi=\alpha_\varphi+c\omega>0,\qquad
 \beta_\varphi^2=(1+c^2)\omega^2.
\]

For a fixed $\omega$ and target $\widehat\theta$, a smooth potential
$\underline\varphi$ is a \emph{$C$-subsolution} if its eigenvalues
$\mu_j=\lambda_j(\alpha+\ddbar\underline\varphi;\omega)$ satisfy
\begin{equation}\label{intro-subsolution-inequality}
 \sum_{j\ne k}\arctan\mu_j>\widehat\theta-\frac\pi2,
 \qquad k=1,\ldots,n,\quad x\in X.
\end{equation}
\begin{theorem}[Lin, preprint {\cite[Theorem 1.5]{Lin23}}]
\label{thm-subsolution-existence}
\SubsolutionStatement
\end{theorem}
This is the formulation in the cited preprint. The published
Collins--Jacob--Yau theorem also assumes
$\vartheta_\omega(\alpha_{\underline\varphi})>(n-2)\pi/2$;
Section~\ref{subsolutions-and-estimates} distinguishes these hypotheses
and gives the estimates \cite[Theorem 1.2]{CJY20}.

For the numerical tests, use
\begin{align}
 \theta&=\frac{n\pi}{2}-\widehat\theta\in(0,\pi),&
 Q_\omega(\eta)&=\frac{n\pi}{2}-\vartheta_\omega(\eta),
 \label{intro-complementary-phase}\\
 \cF_p(A,\Omega;\theta)
 &=\operatorname{Re}(A+\sqrt{-1}\Omega)^p
   -\cot\theta\,\operatorname{Im}(A+\sqrt{-1}\Omega)^p.
 \label{intro-numerical-forms}
\end{align}
Thus $Q_\omega(\eta)=\theta$ is the same equation with the complementary
phase.

\begin{theorem}[Chu--Lee--Takahashi {\cite[Theorem 1.3]{CLT24}}]
\label{thm-ray-criterion}
\RayCriterionStatement
\end{theorem}

\begin{theorem}[Chu--Lee--Takahashi {\cite[Corollary 1.5]{CLT24}}]
\label{thm-projective-criterion}
\ProjectiveCriterionStatement
\end{theorem}
The first criterion is cohomological, so the fixed- and variable-background
existence problems are equivalent for a prescribed pair $(A,\Omega)$.
Projectivity removes the tests with $t>0$; on a general K\"ahler manifold,
endpoint positivity alone can miss the real phase branch
(Example~\ref{ex-torus-phase-lift}).

\subsection{Properness and the choice of scalar flow}
\label{intro-properness-flows}
Fix $\omega$, $\alpha$ and a compatible target $\widehat\theta$.
Define the almost-calibrated potentials and the phase densities by
\begin{align*}
 \cH&=\{\varphi\in C^\infty(X,\R):
 |\vartheta_\omega(\alpha_\varphi)-\widehat\theta|<\pi/2\},\\
 \rho_\varphi&=\operatorname{Re}\bigl(e^{-\sqrt{-1}\widehat\theta}
                   (\omega+\sqrt{-1}\alpha_\varphi)^n\bigr)>0,\\
 \sigma_\varphi&=\operatorname{Im}\bigl(e^{-\sqrt{-1}\widehat\theta}
                   (\omega+\sqrt{-1}\alpha_\varphi)^n\bigr).
\end{align*}
When $0\in\cH$, the functional $\mathcal J$ and the path-length
metric $d_1$ are specified by
\[
 d\mathcal J_\varphi(v)=-\int_Xv\sigma_\varphi,
 \quad\mathcal J(0)=0,
 \qquad \|v\|_{1,\varphi}=\int_X|v|\rho_\varphi.
\]
Here $d_1$ is the infimum of $\int_0^1\|\dot\varphi_t\|_{1,\varphi_t}\,dt$
over smooth paths in $\cH$ joining the two endpoints.
For $0<\theta<\pi/2$, put
\[
 \cH_\varepsilon=\{\varphi\in\cH:
 Q_\omega(\alpha_\varphi)>\varepsilon\},\qquad 0<\varepsilon<\theta.
\]
\begin{theorem}[Chu--Lee {\cite[Theorem 1.4]{CL26}}]
\label{thm-properness}
\PropernessStatement
\end{theorem}
The constants may depend on $\varepsilon$; the statement does not impose
one coercivity constant on all of $\cH$.

A target is \emph{hypercritical} when
$\widehat\theta>(n-1)\pi/2$, equivalently $0<\theta<\pi/2$;
an initial potential is hypercritical when its pointwise phase exceeds
$(n-1)\pi/2$.

For $n\geq2$, with the same fixed data and compatible target, the three
scalar equations have different initial-data conclusions:
\begin{enumerate}
\item If $0<\theta<\pi/2$ and $\varphi_0\in\cH$, then
\[
 \partial_t\varphi=
 \tan\bigl(\vartheta_\omega(\alpha_\varphi)-\widehat\theta\bigr)
\]
exists smoothly for all time. A smooth $C$-subsolution gives smooth
convergence \cite[Theorem 1.4]{Tak20}.
\item If $0<\theta<\pi$ and
$Q_\omega(\alpha_{\varphi_0})<\pi$ on $X$, then
\[
 \partial_t\varphi=\cot Q_\omega(\alpha_\varphi)-\cot\theta
\]
exists smoothly for all time. Convergence follows from a smooth
$C$-subsolution satisfying $Q_\omega(\alpha_{\underline\varphi})<\pi$
\cite[Theorems 1.2--1.3]{FYZ25}.
\item For the original flow
\[
 \partial_t\varphi=\vartheta_\omega(\alpha_\varphi)-\widehat\theta,
\]
if $0<\theta<\pi/2$ and the initial potential is hypercritical
and is itself a $C$-subsolution, the flow exists for all time and
converges smoothly \cite[Theorem 1.3]{Tak20}. Without the
appropriate initial restrictions, finite-time singularities occur
even in smoothly solvable classes \cite{CJ25}.
\end{enumerate}
The distinction matters in
Problem~\ref{problem-original-flow-basin}: its almost-calibrated initial
data need not be hypercritical or a $C$-subsolution.

\subsection{Thomas--Yau: existence, flow and a toric correspondence}
\label{intro-Thomas-Yau}
On a normalized Calabi--Yau manifold $(Y,\omega_Y,\Upsilon)$, a graded
Lagrangian $\Lambda$ has a real phase $\beta_\Lambda$ with
\[
 \omega_Y|_\Lambda=0,\qquad
 \Upsilon|_\Lambda=e^{\sqrt{-1}\beta_\Lambda}\operatorname{vol}_\Lambda.
\]
It is special when $\beta_\Lambda$ is constant. Almost calibrated means
$\operatorname{osc}\beta_\Lambda<\pi$. Write
$Z(\Lambda)=\int_\Lambda\Upsilon=|Z(\Lambda)|e^{\sqrt{-1}\tau(\Lambda)}$
with the phase lift in its phase interval, and normalize
$\tau(\Lambda)=0$. For Lagrangian surfaces, flow stability requires, for every graded
Hamiltonian connected-sum decomposition into almost-calibrated pieces,
$\Lambda\simeq_{\mathrm{Ham}}\Lambda_1\#\Lambda_2$, that
\[
 \begin{gathered}
 [\tau(\Lambda_1),\tau(\Lambda_2)]
 \not\subset(\inf\beta_\Lambda,\sup\beta_\Lambda)
 \quad\text{or}\\
 \operatorname{Area}(\Lambda)
 <|Z(\Lambda_1)|+|Z(\Lambda_2)|.
 \end{gathered}
\]
The connected sum is Lagrangian neck surgery with gradings extending
over the necks. 
The Thomas--Yau
existence question seeks phase-stability conditions characterizing a
special representative in a Hamiltonian isotopy class
\cite[Section 1]{TY02}. It must be separated from a smooth-flow assertion:
\begin{itemize}
\item Neves proves that, on a compact Calabi--Yau surface, every compact
embedded Lagrangian has a Hamiltonian-isotopic embedded representative
whose mean curvature flow becomes singular in finite time. This applies
even when the starting class contains a special Lagrangian. His relevant
examples are not almost calibrated \cite[Theorem 6.1 and Remark 1.3]{Nev13}.
\item Lotay--Oliveira prove smooth long-time existence and convergence
to a special Lagrangian for compact, embedded, circle-invariant,
almost-calibrated and flow-stable Lagrangians in Gibbons--Hawking
ALE/ALF hyperk\"ahler four-manifolds \cite[Theorem 1.2]{LO24}.
\end{itemize}
Section~\ref{Thomas-Yau} distinguishes this condition
from the historical Thomas--Yau quantifiers and from Joyce's programme
of categorical representatives and flow with surgery \cite{Joy15}.

A toric manifold is \emph{weak Fano} when $-K_X$ is nef and big.
For the toric correspondence, let $X$ be a smooth projective toric
weak Fano $n$-fold, $\omega$ a torus-invariant K\"ahler form, and $M$
a line bundle with ample dual. Let $\mathcal L_k$ denote the SYZ transform
of $M^{\otimes k}$ over the interior of the moment polytope of $(X,k\omega)$:
it is the Lagrangian section determined by a Hermitian metric on that
line bundle. Write $(\mathcal Y_{q_k},W_k)$ for the toric
Landau--Ginzburg mirror, and set
\[
 \mathcal Y_{q_k}\cong(\C^*)^n,\qquad
 \Upsilon_0=\bigwedge_{j=1}^n\frac{dz^j}{z^j},\qquad
 P_k(\mathcal E)=\int_{\gamma(\mathcal E)}e^{-W_k}\Upsilon_0.
\]
Here $\gamma(\mathcal E)$ is the rapid-decay cycle of a
Fukaya--Seidel object; its ends lie where $\operatorname{Re}W_k\to+\infty$.
A shift by $r$ multiplies this period by $(-1)^r$.
\Needspace{17\baselineskip}
\begin{theorem}[Stoppa, preprint {\cite[Theorem 1.2 and Corollary 2.7]{Sto25}}]
\label{thm-toric-mirror}
\ToricMirrorStatement
\end{theorem}
The statement is from the cited v2 preprint. Special SYZ sections
use the geometric mirror volume form, not the oscillatory form defining
$P_k$. The theorem uses one specified family of test objects, not every
subobject in an arbitrary heart. In this setting, its period test is
equivalent to smooth supercritical dHYM solvability on $M^\vee$;
Section~\ref{toric-mirror-correspondence} gives the signs and scaling
\cite[Section 2]{Sto25}.

\subsection{Surface stability does not imply smooth solvability}
\label{intro-surface-stability}
On a smooth projective surface and for an ample
$\Xi\in\NS(X)_\R$, the divisorial stability condition $\sigma_{0,\Xi}$
has central charge
\[
 Z_{0,\Xi}(E)=-\big\langle e^{-\sqrt{-1}\Xi}\operatorname{ch}(E),[X]\big\rangle
\]
and heart obtained by tilting $\operatorname{Coh}(X)$ at slope zero.
More explicitly, let $\mathcal T_{0,\Xi}$ be generated by torsion sheaves
and slope-semistable torsion-free sheaves of positive slope, and
$\mathcal F_{0,\Xi}$ by slope-semistable torsion-free sheaves of
nonpositive slope, where
$\mu_\Xi(F)=\langle c_1(F)\smile\Xi,[X]\rangle/\operatorname{rk}(F)$.
The heart is the extension closure
$\mathcal A_{0,\Xi}=\langle\mathcal F_{0,\Xi}[1],\mathcal T_{0,\Xi}\rangle$.
Stability means
\[
 \operatorname{Arg}Z_{0,\Xi}(G)<\operatorname{Arg}Z_{0,\Xi}(E)
 \quad(0\ne G\subsetneq E\text{ in }\mathcal A_{0,\Xi}),
 \qquad\operatorname{Arg}\in(0,\pi].
\]
\begin{theorem}[M\"akel\"a, preprint {\cite[Corollary 1.5]{Mak26}}]
\label{thm-stability-under-scaling}
\SurfaceStabilityStatement
\end{theorem}
This is the statement of the August 2026 preprint \cite{Mak26}.
The non-strict curve inequalities differ from the strict inequalities
in Theorem~\ref{thm-surface-criterion}. In particular, on
\[
 X=\operatorname{Bl}_p\bbP^2,\qquad
 \Omega=\frac{2H-E}{\sqrt3},\qquad L=\mathcal O_X(H),
\]
where $H$ is the pullback hyperplane class and $E$ is the exceptional
divisor, one has $A_c=H$, which is nef but not K\"ahler. Consequently,
\[
 \left.
 \begin{array}{c}
 L^{\otimes k}\text{ is }\sigma_{0,k\Omega}\text{-stable}\\
 \text{for every integer }k\geq1
 \end{array}\right\}\quad\text{but}\quad
 \text{no smooth dHYM solution exists}.
\]
The calculation is given in Section~\ref{surface-stability}
\cite[Example 2.11]{Fan26}.

\Needspace{10\baselineskip}
At a fixed scale, the same preprint \cite[Theorem 1.1(a)]{Mak26}
gives the following detector. Let $X$ be a smooth projective surface,
$\Omega\in\NS(X)_\R$ ample, and $L$ a line bundle with
$s=\Omega\cdot c_1(L)>0$. Then
\begin{equation}\label{intro-divisor-detector}
 \begin{gathered}
 L\text{ is not }\sigma_{0,\Omega}\text{-stable}
 \quad\Longleftrightarrow\quad
 \exists\,D\ne0\text{ effective with}\\
 D^2<0,\qquad 0<\Omega\cdot D<s,\qquad
 \operatorname{Arg}Z_{0,\Omega}(L(-D))
 \geq\operatorname{Arg}Z_{0,\Omega}(L).
 \end{gathered}
\end{equation}
Here $D$ need not be reduced or irreducible. The full $B$-twisted
statement and the stronger integral-curve question are separated in
Theorem~\ref{thm-effective-divisor-detector} and
Problem~\ref{problem-integral-curve-detector}.


\section{The Thomas--Yau problem}\label{Thomas-Yau}
\subsection{Phase, calibration and existence}
Let $(Y,\omega_Y,J_Y,\Upsilon)$ be a Calabi--Yau $n$-fold with
Ricci-flat metric $g_Y=\omega_Y(\,\cdot\,,J_Y\,\cdot\,)$ and parallel
holomorphic volume form $\Upsilon$, normalized so that its restriction
to any oriented Lagrangian plane has unit norm. A compact oriented
Lagrangian $\Lambda\subset Y$ satisfies
\[
 \omega_Y|_\Lambda=0,\qquad
 \Upsilon|_\Lambda=e^{\sqrt{-1}\beta_\Lambda}
       \operatorname{vol}_\Lambda.
\]
A \emph{grading} is a continuous real lift $\beta_\Lambda$ of this phase;
its existence is the zero-Maslov condition. The Lagrangian is
\emph{special} when the lift is constant. It is \emph{almost calibrated}
when
\[
 \sup_\Lambda\beta_\Lambda-\inf_\Lambda\beta_\Lambda<\pi.
\]
In that case
\[
 Z(\Lambda):=\int_\Lambda\Upsilon\ne0,\qquad
 Z(\Lambda)=|Z(\Lambda)|e^{\sqrt{-1}\tau(\Lambda)},\qquad
 \tau(\Lambda)\in[\inf_\Lambda\beta_\Lambda,
                    \sup_\Lambda\beta_\Lambda].
\]
This specifies the real lift of the cohomological phase
\cite[Sections 1--2]{TY02}.

Write $\Lambda'\simeq_{\mathrm{Ham}}\Lambda$ if they are connected by
an isotopy $F_s$ for which
$F_s^*(\iota_{\partial_sF_s}\omega_Y)$ is exact.
The period $Z$ is constant along such an isotopy. If
$Z(\Lambda)=|Z(\Lambda)|e^{\sqrt{-1}\tau}$, then
\begin{align}
 \operatorname{Vol}(\Lambda')-|Z(\Lambda)|
 &=\int_{\Lambda'}
   \bigl(1-\cos(\beta_{\Lambda'}-\tau)\bigr)
   \operatorname{vol}_{\Lambda'}\geq0.\label{calibration-deficit}
\end{align}
Equality holds precisely when $\Lambda'$ is calibrated by
$\operatorname{Re}(e^{-\sqrt{-1}\tau}\Upsilon)$.
This lower bound does not supply compactness of smooth minimizers.

\begin{problem}[Thomas--Yau {\cite[Section 1]{TY02}}]
\label{problem-Thomas-Yau-existence}
Let $(Y,\omega_Y,J_Y,\Upsilon)$ be a normalized Calabi--Yau $n$-fold
and let $\Lambda\subset Y$ be a compact graded Lagrangian with
$Z(\Lambda)\ne0$. Choose $\tau\in\R$ with
$e^{-\sqrt{-1}\tau}Z(\Lambda)>0$. Determine phase-stability conditions on the allowed
decompositions of $\Lambda$ that characterize the existence of a smooth
representative $\Lambda_*\simeq_{\mathrm{Ham}}\Lambda$ satisfying
\[
 \Upsilon|_{\Lambda_*}
 =e^{\sqrt{-1}\tau}\operatorname{vol}_{\Lambda_*}.
\]
\end{problem}
This is the existence question underlying the Thomas--Yau programme,
not a definition of a universal stability condition. The choice of
allowed decompositions, and whether one seeks a smooth representative
or a singular categorical representative, are part of its formulation.

For the local graphical calculation, take
$Y=U\times\R^n\subset\C^n$, with $z^j=x^j+\sqrt{-1}y^j$,
$\omega_Y=\sum_j dx^j\wedge dy^j$ and
$\Upsilon=\bigwedge_j dz^j$. For $u\in C^\infty(U,\R)$, the graph
$F_u(x)=(x,Du(x))$ satisfies
\begin{align}
 F_u^*\omega_Y
 &=\sum_{j,k}u_{jk}\,dx^j\wedge dx^k=0,\nonumber\\
 F_u^*\Upsilon
 &=\det(I+\sqrt{-1}D^2u)\,dx^1\wedge\cdots\wedge dx^n,\nonumber\\
 \beta_{F_u}&=\sum_{j=1}^n\arctan\lambda_j(D^2u).
 \label{graph-phase-calculation}
\end{align}
The scalar dHYM operator replaces this real Hessian by a Hermitian
$(1,1)$-form. The common determinant explains the phase analogy;
a global mirror correspondence requires additional geometry.

\subsection{Initial conditions and mean curvature flow}
For a smooth Lagrangian mean curvature flow
$\partial_tF=H$, with mean curvature vector $H$, one has
\begin{equation}\label{Lagrangian-flow-identities}
 H=J_Y\nabla\beta,\qquad
 \partial_t\beta=\Delta_{\Lambda_t}\beta,\qquad
 \frac{d}{dt}\operatorname{Vol}(\Lambda_t)
 =-\int_{\Lambda_t}|H|^2\operatorname{vol}_{\Lambda_t},
\end{equation}
where $\Delta=\operatorname{div}\nabla$ \cite[Section 2]{TY02}.
Thus phase extrema move inward and volume decreases for as long as the
flow is smooth. Neither identity controls curvature by itself.

A \emph{graded connected sum} $\Lambda_1\#\Lambda_2$ is an ordered
Lagrangian neck surgery at intersection points, with the gradings of
the two pieces extending over the necks. For the historical
Thomas--Yau conditions, the allowed decompositions include the relative
versions in which surgery is performed in a family
\cite[Sections 3 and 7]{TY02}. With the phase lifts of that construction,
write, for a decomposition $D=(\Lambda_1,\Lambda_2)$ of $\Lambda$,
\begin{align}
 \mathsf A_D:
 &\quad [\tau(\Lambda_1),\tau(\Lambda_2)]
       \not\subset
       (\inf_\Lambda\beta_\Lambda,\sup_\Lambda\beta_\Lambda),
       \label{TY-phase-test}\\
 \mathsf B_D^{\leq}:
 &\quad \operatorname{Vol}(\Lambda)
       \leq |Z(\Lambda_1)|+|Z(\Lambda_2)|.
       \label{TY-volume-test}
\end{align}
After normalizing $\tau(\Lambda)=0$, the original Conjecture~7.3 of
\cite{TY02} predicts smooth long-time existence and convergence to a
special representative under
\[
 (\forall D,\ \mathsf A_D)
 \quad\text{or}\quad
 (\forall D,\ \mathsf B_D^{\leq}).
\]
These are restrictions on the initial representative, not merely the
assertion that its Hamiltonian class contains a special Lagrangian.
The displayed conditions record the original formulation; the theorem
below uses a precise almost-calibrated variant.

\begin{theorem}[Neves {\cite[Theorem 6.1]{Nev13}}]
\label{thm-Neves}
Let $Y$ be a compact Calabi--Yau surface and let $\Sigma\subset Y$ be a compact
embedded Lagrangian. There exists a compact embedded Lagrangian
$\Lambda\simeq_{\mathrm{Ham}}\Sigma$ whose mean curvature flow develops
a singularity at a finite time.
\end{theorem}
Taking $\Sigma$ special Lagrangian shows that existence of a calibrated
representative is insufficient for arbitrary initial data. These
examples are not almost calibrated
\cite[Remark 1.3(4)]{Nev13}; they do not refute the flow statement with
its additional initial conditions.

\begin{definition}\label{def-flow-stability}
Let $\Lambda$ be a compact almost-calibrated Lagrangian in a
Calabi--Yau surface. Rotate $\Upsilon$ so that $\tau(\Lambda)=0$.
It is \emph{flow stable} in the sense of
\cite[Definition 6.2]{LO24} if, for every graded decomposition
$\Lambda\simeq_{\mathrm{Ham}}\Lambda_1\#\Lambda_2$ into
almost-calibrated Lagrangians, at least one of the following holds:
\begin{itemize}
\item
$[\tau(\Lambda_1),\tau(\Lambda_2)]
 \not\subset(\inf_\Lambda\beta_\Lambda,\sup_\Lambda\beta_\Lambda)$;
\item
$\operatorname{Area}(\Lambda)
 <|Z(\Lambda_1)|+|Z(\Lambda_2)|$.
\end{itemize}
\end{definition}
The disjunction is imposed separately on each decomposition, and the
area inequality is strict. Neither convention should be replaced by
the quantifiers or weak inequality of the original formulation.

\begin{theorem}[Lotay--Oliveira {\cite[Theorem 1.2]{LO24}}]
\label{thm-LO}
Let $Y$ be an ALE or ALF hyperk\"ahler four-manifold obtained from the
Gibbons--Hawking construction. Fix one of its compatible Calabi--Yau
structures. Suppose $\Lambda\subset Y$ is
\begin{itemize}
\item compact, embedded and invariant under the Gibbons--Hawking circle action;
\item almost calibrated and flow stable in the sense of
Definition~\ref{def-flow-stability}.
\end{itemize}
Then its Lagrangian mean curvature flow exists smoothly for all
$t\geq0$ and converges smoothly to a circle-invariant special Lagrangian
in its Hamiltonian class.
\end{theorem}
This is a result in a specified noncompact ambient geometry with
symmetry, not a theorem for arbitrary Calabi--Yau surfaces.

\subsection{Categorical representatives and surgery}
Joyce's formulation uses an enlarged derived Fukaya category with
objects represented by branes $(\Lambda,E,b)$: graded Lagrangians,
rank-one local systems $E$, and Floer bounding cochains $b$ removing
obstructions. Its proposed central charge is
\[
 Z([\Lambda,E,b])=\int_\Lambda\Upsilon.
\]
The expected relation is between semistability and constant-phase
representatives in an appropriately enlarged class, which can include
immersed or singular Lagrangians
\cite[Conjecture 3.2]{Joy15}.

The flow programme starts with generic Floer-unobstructed branes and
allows surgeries. At regular times, it preserves the object up to
isomorphism in the category; the proposed limit is assembled from
special Lagrangian integral currents with ordered phases
\cite[Conjecture 3.34]{Joy15}. It does not assert that topology is
preserved through surgery, or that every limit is a smooth embedding.
Hamiltonian isotopy, categorical isomorphism and equality of homology
classes are different equivalence relations. The toric result in
Section~\ref{toric-mirror-correspondence} concerns special SYZ sections;
the scalar flow question in Section~\ref{properness-and-evolution}
concerns smooth potentials. Neither is the entire surgery programme.

\section{The surface criterion}\label{surface-criterion}
Let $X$ be a compact K\"ahler surface. Set
\[
 v=\langle\Omega^2,[X]\rangle,\qquad
 a=\langle A^2,[X]\rangle,\qquad
 s=\langle A\smile\Omega,[X]\rangle.
\]
Assume $s>0$. Then $Z_X(A,\Omega)=v-a+2\sqrt{-1}s$ lies in the upper
half-plane. Set
\begin{equation}\label{shifted-surface-class}
 c=\frac{v-a}{2s},\qquad
 \widehat\theta=\operatorname{arccot}c\in(0,\pi),\qquad A_c=A+c\Omega.
\end{equation}

\begin{theoremrecall}{thm-surface-criterion}
\SurfaceCriterionStatement
\end{theoremrecall}

\begin{proof}
Expanding the phase equation and completing the square gives
\begin{align}
 \alpha_\varphi^2+2c\,\alpha_\varphi\wedge\omega&=\omega^2,\nonumber\\
 \beta_\varphi:=\alpha_\varphi+c\omega,\qquad
 [\beta_\varphi]&=A_c,\qquad
 \beta_\varphi^2=(1+c^2)\omega^2.
 \label{surface-Monge-Ampere-equation}
\end{align}
At a positive-phase solution, the eigenvalues of $\alpha_\varphi$ satisfy
\[
 \lambda_j+c=\frac{1+\lambda_j^2}{\lambda_1+\lambda_2}>0,
 \qquad j=1,2.
\]
Thus $\beta_\varphi>0$, proving \textup{(i)}$\Rightarrow$\textup{(ii)}.

Conversely, suppose $A_c\in\cK_X$ and choose any K\"ahler form
$\widetilde\omega\in\Omega$. The identity
$\langle A_c^2,[X]\rangle=(1+c^2)v$ and Yau's theorem \cite{Yau78}
give a K\"ahler form $\widetilde\beta\in A_c$ such that
\[
 \widetilde\beta^2=(1+c^2)\widetilde\omega^2.
\]
Set $\widetilde\eta=\widetilde\beta-c\widetilde\omega\in A$.
Its eigenvalues $\lambda_1,\lambda_2$ relative to $\widetilde\omega$ satisfy
\[
 (\lambda_1+c)(\lambda_2+c)=1+c^2,\qquad
 \lambda_1+\lambda_2\geq2\sqrt{1+c^2}-2c>0.
\]
Thus $\vartheta_{\widetilde\omega}(\widetilde\eta)\in(0,\pi)$ and
$\cot\vartheta_{\widetilde\omega}(\widetilde\eta)=c$; hence
$\vartheta_{\widetilde\omega}(\widetilde\eta)=\widehat\theta$.

By the Hodge index theorem, $av\leq s^2$. Hence
\[
 \langle A_c^2,[X]\rangle=(1+c^2)v>0,\qquad
 \langle A_c\smile\Omega,[X]\rangle=s+cv
 \geq\frac{s^2+v^2}{2s}>0.
\]
Thus $A_c$ lies in the positive-cone component containing $\Omega$.
The surface K\"ahler-cone criterion \cite{DP04} gives
\textup{(ii)}$\Leftrightarrow$\textup{(iii)}.
\end{proof}
For fixed $\omega$, the Monge--Amp\`ere equation also gives uniqueness
of the potential modulo constants.

A target is \emph{hypercritical} when
$(n-1)\pi/2<\widehat\theta<n\pi/2$. If the K\"ahler class is not
prescribed, there is an elementary construction in every dimension:
\[
 \eta\in A,\quad\eta>0,\qquad
 \omega=\frac{\eta}{\tan(\widehat\theta/n)}
 \quad\Longrightarrow\quad\vartheta_\omega(\eta)=\widehat\theta.
\]
Conversely, a hypercritical solution has $\eta>0$, since a nonpositive
eigenvalue would force the phase below $(n-1)\pi/2$.
This gives existence for some freely chosen background exactly when
$A$ is K\"ahler. It changes $\Omega$ and is not a criterion for a
prescribed pair $(A,\Omega)$.


\begin{example}\label{ex-nef-boundary-class}
On $X=\operatorname{Bl}_p\bbP^2$, let $H$ be the pullback hyperplane
class and $E$ the exceptional divisor class. Take
\[
 \Omega=\frac{2H-E}{\sqrt3},\qquad A=H.
\]
Then $v=a=1$, $c=0$ and $A_c=H$. This class is nef and big, but
$\langle A_c,[E]\rangle=0$. Thus $A_c\notin\cK_X$, and no K\"ahler
background $\omega\in\Omega$ admits a smooth dHYM solution.
\end{example}

\section{Subsolutions and estimates}\label{subsolutions-and-estimates}
For a holomorphic line bundle $L$ with Hermitian metric $h$, the Chern
curvature in a local holomorphic frame $e$ is
\[
 F_h=-\partial\overline\partial\log h(e,e),\qquad
 c_1(L,h)=\frac{\sqrt{-1}}{2\pi}F_h.
\]
If $A=c_1(L)$ and $\alpha=c_1(L,h_0)$, then
\[
 h=e^{-2\pi\varphi}h_0
 \quad\Longrightarrow\quad c_1(L,h)=\alpha+\ddbar\varphi.
\]
Thus the potential formulation includes the metric equation on a line
bundle \cite{JY17}.


For a variation $v\in C^\infty(X,\R)$, choose coordinates unitary for
$\omega$ and diagonalizing $\alpha_\varphi$. Then
\[
 D_\varphi\bigl(\vartheta_\omega(\alpha_\varphi)\bigr)[v]
 =\sum_{j=1}^n\frac{v_{j\overline j}}{1+\lambda_j^2}.
\]
The linearization is elliptic, but its ellipticity constants degenerate
when $|\lambda_j|\to\infty$.

\begin{definition}\label{def-smooth-subsolution}
Fix a K\"ahler background $\omega$ and a target phase
$\widehat\theta\in((n-2)\pi/2,n\pi/2)$.
A potential $\underline\varphi\in C^\infty(X,\R)$ is a
$C$-subsolution for $\vartheta_\omega(\alpha_\varphi)=\widehat\theta$
if the eigenvalues $\mu_j$ of
$\alpha_{\underline\varphi}=\alpha+\ddbar\underline\varphi$
relative to $\omega$ satisfy
\[
 \sum_{j\ne k}\arctan\mu_j>\widehat\theta-\frac\pi2,
 \qquad k=1,\ldots,n,\quad x\in X.
\]
\end{definition}

For fixed $x\in X$, the set
\[
 \left\{\lambda\in\R^n:\lambda_j\geq\mu_j(x)\ (j=1,\ldots,n),\quad
 \sum_{j=1}^n\arctan\lambda_j=\widehat\theta\right\}
\]
is bounded. Indeed, if $\lambda_k\to+\infty$, then
\[
 \liminf\sum_j\arctan\lambda_j
 \geq\frac\pi2+\sum_{j\ne k}\arctan\mu_j(x)>\widehat\theta,
\]
a contradiction \cite[Lemma 3.3]{CJY20}.
On a surface,
\[
 n=2:\qquad \mu_j+\cot\widehat\theta>0\ (j=1,2)
 \quad\Longleftrightarrow\quad
 \alpha_{\underline\varphi}+\cot\widehat\theta\,\omega>0.
\]

\begin{theoremrecall}{thm-subsolution-existence}
\SubsolutionStatement
\end{theoremrecall}

The published Collins--Jacob--Yau theorem assumes in addition
\[
 \vartheta_\omega(\alpha_{\underline\varphi})>\frac{(n-2)\pi}{2}
 \qquad\text{\cite[Theorem 1.2]{CJY20}}.
\]
Summing the $C$-subsolution inequalities gives
\[
 (n-1)\vartheta_\omega(\alpha_{\underline\varphi})
 >n\left(\widehat\theta-\frac\pi2\right).
\]
Thus the extra assumption is automatic when
$\widehat\theta>(n-2+2/n)\pi/2$, but not throughout the supercritical range.

Under the Collins--Jacob--Yau hypotheses, the continuity argument
uses a smoothed maximum of the initial and target phase functions,
then deforms to the constant $\widehat\theta$
\cite[Section 7]{CJY20}. For its solution potentials, normalized by
$\sup_X\varphi=0$, the estimates include
\begin{align*}
 \|\varphi\|_{C^0(X)}&\leq C,\\
 \sup_X|\ddbar\varphi|_\omega
 &\leq C\bigl(1+\sup_X|\nabla\varphi|_\omega^2\bigr)
 \qquad\text{\cite[Proposition 3.6 and Theorem 4.1]{CJY20}}.
\end{align*}
If $M_j=\sup_X|\nabla\varphi_j|_\omega\to\infty$, rescaling at points
where the maxima occur produces a limit $u$ satisfying
\[
 u\colon\C^n\longrightarrow\R,\qquad
 \|u\|_{L^\infty}<\infty,\qquad
 \sqrt{-1}\partial\overline\partial u\geq0,\qquad |\nabla u(0)|=1.
\]
The restriction of $u$ to each complex line is bounded and subharmonic,
hence constant. This contradicts $|\nabla u(0)|=1$
\cite[Proposition 5.1]{CJY20}. The gradient bound then gives the Hessian
bound, uniform ellipticity and higher regularity.

\section{Numerical criteria}\label{intersection-criteria}
Use $\operatorname{arccot}\lambda\in(0,\pi)$ and set
\begin{equation}\label{complementary-phase}
 \theta=\frac{n\pi}{2}-\widehat\theta\in(0,\pi),\qquad
 Q_\omega(\eta)=\sum_{j=1}^n\operatorname{arccot}\lambda_j(\eta;\omega)
 =\frac{n\pi}{2}-\vartheta_\omega(\eta).
\end{equation}
Then the dHYM equation is $Q_\omega(\alpha_\varphi)=\theta$. For a
closed real $(1,1)$-form $\eta$, define
\[
 F_p(\eta,\omega;\theta)
 =\operatorname{Re}(\eta+\sqrt{-1}\omega)^p
 -\cot\theta\,\operatorname{Im}(\eta+\sqrt{-1}\omega)^p.
\]
At a solution, $F_n(\alpha_\varphi,\omega;\theta)=0$.
For a complex $p$-plane $\Pi\subset T_xX$, $1\leq p<n$, let
$\nu_1\leq\cdots\leq\nu_p$ be the eigenvalues of
$\alpha_\varphi|_\Pi$ relative to $\omega|_\Pi$. Set
\[
 r_\Pi=\prod_{j=1}^p\sqrt{1+\nu_j^2},\qquad
 q_\Pi=\sum_{j=1}^p\operatorname{arccot}\nu_j.
\]
Ordering the full eigenvalues as $\lambda_1\leq\cdots\leq\lambda_n$,
eigenvalue interlacing gives
\[
 \nu_j\geq\lambda_j,\qquad
 0<q_\Pi\leq\sum_{j=1}^p\operatorname{arccot}\lambda_j
 <Q_\omega(\alpha_\varphi)=\theta.
\]
Hence
\begin{align*}
 \left.(\alpha_\varphi+\sqrt{-1}\omega)^p\right|_\Pi
 &=r_\Pi e^{\sqrt{-1}q_\Pi}\left.\omega^p\right|_\Pi,\\
 \left.F_p(\alpha_\varphi,\omega;\theta)\right|_\Pi
 &=r_\Pi\frac{\sin(\theta-q_\Pi)}{\sin\theta}
       \left.\omega^p\right|_\Pi>0.
\end{align*}
Define the cohomology classes
\begin{align}
 \cF_p(A,\Omega;\theta)
 &:=\operatorname{Re}(A+\sqrt{-1}\Omega)^p
       -\cot\theta\,\operatorname{Im}(A+\sqrt{-1}\Omega)^p\nonumber\\
 &=[F_p(\alpha_\varphi,\omega;\theta)]\in H^{p,p}(X,\R),
 \qquad \cF_0=1.
 \label{numerical-cohomology-class}
\end{align}
Integration gives the necessary inequalities
\[
 \langle\cF_p(A,\Omega;\theta),[V]\rangle>0,
 \qquad 1\leq p=\dim V<n,
\]
for every irreducible analytic subvariety $V\subset X$
\cite[Section 8]{CJY20}. On a surface,
\[
 \theta=\pi-\widehat\theta,\qquad
 \cF_1(A,\Omega;\theta)=A+\cot(\widehat\theta)\Omega=A_c.
\]

\begin{example}
\label{ex-torus-phase-lift}
Let $X$ be a complex three-torus with no proper positive-dimensional
analytic subvarieties. Take $\alpha=-\omega$, so $A=-\Omega$.
The target $\widehat\theta=5\pi/4$, or $\theta=\pi/4$, satisfies
\eqref{phase-compatibility}; all proper-subvariety tests are vacuous.
At a maximum of $\varphi$, however,
\[
 \alpha_\varphi\leq-\omega
 \quad\Longrightarrow\quad
 \vartheta_\omega(\alpha_\varphi)\leq-\frac{3\pi}{4}
 \neq\frac{5\pi}{4}.
\]
Thus the equation with target $\widehat\theta=5\pi/4$ has no solution
\cite{Zha23}.
\end{example}

The endpoint tests do not detect this failure. Along the ray
$A+t\Omega$,
\[
 \frac{\langle\cF_3(A+t\Omega,\Omega;\pi/4),[X]\rangle}
      {\langle\Omega^3,[X]\rangle}
 =t^3-6t^2+6t,
\]
and the value at $t=2$ is $-4$. The criterion below requires the
corresponding inequalities for every $t\geq0$.

\Needspace{13\baselineskip}
\begin{theoremrecall}{thm-ray-criterion}
\RayCriterionStatement
\end{theoremrecall}

Consequently, for compact K\"ahler $X$, solvability for one
$\omega\in\Omega$ is equivalent to solvability for every such
representative: all the conditions of Theorem~\ref{thm-ray-criterion}
are cohomological.

The proof in \cite[Sections 3--7]{CLT24} constructs local potentials
$u_1,\ldots,u_m$ by mass concentration and regularization on resolutions,
then glues them using regularized maxima. On an overlap, write
$u=(u_1,\ldots,u_m)$. For a smooth convex $M\colon\R^m\to\R$ with
$\partial_iM\geq0$ and $\sum_i\partial_iM=1$,
\begin{align*}
 \alpha+\ddbar M(u)
 &=\sum_i\partial_iM(u)\,(\alpha+\ddbar u_i)\\
 &{}+\sqrt{-1}\sum_{i,j}\partial_i\partial_jM(u)\,
          \partial u_i\wedge\overline\partial u_j.
\end{align*}
The first term is a convex combination of the local forms, and the
second is semipositive. The admissible set used in
\cite[Sections 3--7]{CLT24} is convex and is preserved by adding a
semipositive $(1,1)$-form; hence the glued form remains admissible.

\begin{theoremrecall}{thm-projective-criterion}
\ProjectiveCriterionStatement
\end{theoremrecall}

The additional implication in Theorem~\ref{thm-projective-criterion} is
\[
 \left.
 \begin{array}{c}
  \langle\cF_p(A,\Omega;\theta),[V]\rangle>0,
  \quad 1\leq p=\dim V<n,\\[2pt]
  \langle\cF_n(A,\Omega;\theta),[X]\rangle=0
 \end{array}
 \right\}
 \Longrightarrow
 \begin{gathered}\text{the ray inequalities}\\\text{of Theorem~\ref{thm-ray-criterion}}\end{gathered}.
\]
For a very ample integral class $\Gamma$, the relevant intersection
number is the polynomial
\begin{align}
 &\big\langle\cF_p(A+t\Gamma,\Omega;\theta),[V]\big\rangle\nonumber\\
 &=\sum_{k=0}^p\binom pk t^{p-k}
 \big\langle\cF_k(A,\Omega;\theta)\smile\Gamma^{p-k},[V]\big\rangle.
 \label{hyperplane-section-expansion}
\end{align}
For $0<k<p$, a generic $(p-k)$-fold hyperplane section represents
$\Gamma^{p-k}\cap[V]$ by an effective $k$-cycle
$\sum_\ell m_\ell[W_\ell]$. Thus
\[
 \big\langle\cF_k(A,\Omega;\theta)\smile\Gamma^{p-k},[V]\big\rangle
 =\sum_\ell m_\ell\langle\cF_k(A,\Omega;\theta),[W_\ell]\rangle>0.
\]
The coefficient with $k=0$ is $\langle\Gamma^p,[V]\rangle>0$.
The constant term is positive for $p<n$ and zero for $p=n$, by the
hypotheses of Theorem~\ref{thm-projective-criterion}.
The test-family theorem of \cite{CLT24} therefore applies to
$A+t\Gamma$.

\section{Toric mirror correspondence}\label{toric-mirror-correspondence}
Let $X$ be a smooth projective toric $n$-fold with $-K_X$ nef and big
(a \emph{toric weak Fano} manifold). Fix a torus-invariant K\"ahler
form $\omega\in\Omega$ and a line bundle $M$ with ample dual. Use
\[
 A=c_1(M^\vee)=-c_1(M),\qquad
 \theta=\frac{n\pi}{2}-\widehat\theta\in(0,\pi),\qquad
 e^{-\sqrt{-1}\widehat\theta}Z_X(A,\Omega)>0.
\]
The dual is essential: in the toric SYZ convention used in
\cite[Remark 2.1]{Sto25}, the special-section equation for $M$ is the
positive-sign dHYM equation \eqref{dHYM-equation} for $M^\vee$.

For a toric subvariety $V$ of dimension $p$, put
\[
 T_V(M)=\int_V e^{-\sqrt{-1}\Omega}\operatorname{ch}(M),\qquad
 W_p(V)=\langle(A+\sqrt{-1}\Omega)^p,[V]\rangle.
\]
Since $\operatorname{ch}(M)=e^{-A}$,
\begin{align}
 T_V(M)&=\frac{(-1)^p}{p!}W_p(V),&
 W_n(X)&=R e^{\sqrt{-1}\theta},\quad R>0,\nonumber\\
 \operatorname{Im}\frac{(-1)^{n-p}T_V(M)}{T_X(M)}
 &=-\frac{n!\sin\theta}{p!R}
   \langle\cF_p(A,\Omega;\theta),[V]\rangle.
 \label{numerical-to-phase-ratio}
\end{align}
Thus the numerical inequality is a strict phase inequality, with the
codimension sign already fixed. For nonzero $z,w$, the convention
\[
 \operatorname{Im}(z/w)<0
 \quad\Longleftrightarrow\quad
 \operatorname{Arg}_0(z/w)\in(-\pi,0),
 \qquad\operatorname{Arg}_0\in(-\pi,\pi],
\]
avoids independently choosing arguments modulo $2\pi$.

For $k>0$, the analytic equation and its numerical tests are exactly scale invariant:
\begin{align}
 \vartheta_{k\omega}(k\eta)&=\vartheta_\omega(\eta),\nonumber\\
 \cF_p(kA,k\Omega;\theta)&=k^p\cF_p(A,\Omega;\theta).
 \label{toric-simultaneous-scaling}
\end{align}

For the Landau--Ginzburg mirror $(\mathcal Y_{q_k},W_k)$ of
$(X,k\omega)$, let
\[
 \mathcal Y_{q_k}\cong(\C^*)^n,\qquad
 \Upsilon_0=\frac{dz^1}{z^1}\wedge\cdots\wedge\frac{dz^n}{z^n},
 \qquad
 P_k(\mathcal E)=\int_{\gamma(\mathcal E)}e^{-W_k}\Upsilon_0.
\]
Here $\gamma(\mathcal E)$ is the rapid-decay integration cycle of an
object in the derived Fukaya--Seidel category: its noncompact ends lie
where $\operatorname{Re}W_k\to+\infty$. A shift by $r$ multiplies its
period by $(-1)^r$. The SYZ transform of a Hermitian line bundle is a
Lagrangian section over the interior of the moment polytope; changing
the metric preserves its Hamiltonian class
\cite[Section 1.1]{Sto25}.

\Needspace{18\baselineskip}
\begin{theoremrecall}{thm-toric-mirror}
\ToricMirrorStatement
\end{theoremrecall}
The nonvanishing assumption is the genericity condition of
\cite[Remark 2.6(v)]{Sto25}. The special-section equation uses the
mirror geometric volume form; the oscillatory form
$e^{-W_k}\Upsilon_0$ defines the periods in the test of
Theorem~\ref{thm-toric-mirror}. These roles should not be conflated.

The mirror comparison instead uses large-volume asymptotics of
Iritani's toric $\widehat\Gamma$-theorem. There are only finitely many
toric tests, and their nonzero leading terms let one choose a common
scale preserving every sign \cite[Theorem 2.5 and Section 5.2]{Sto25}.
Together with \eqref{numerical-to-phase-ratio} and
Theorem~\ref{thm-projective-criterion}, this gives the sequence of
\emph{proved equivalences in this setting}
\[
 \begin{array}{c}
 \text{strict numerical inequalities for }(A,\Omega)\\
 \Updownarrow\\
 \text{smooth supercritical dHYM on }M^\vee\\
 \Updownarrow\\
 \text{special SYZ representative of }\mathcal L_k
 \quad\Longleftrightarrow\quad\text{the period test of Theorem~\ref{thm-toric-mirror}}.
 \end{array}
\]
The test objects are not asserted to exhaust all subobjects in an
arbitrary categorical heart. A general Bridgeland-stability equivalence
would be an additional statement.

\section{Properness and flows}\label{properness-and-evolution}
Fix a K\"ahler form $\omega\in\Omega$ and a target $\widehat\theta$
satisfying \eqref{phase-compatibility}. Relative to this fixed target,
the almost-calibrated potentials form the space
\[
 \cH=\{\varphi\in C^\infty(X,\R):
 |\vartheta_\omega(\alpha_\varphi)-\widehat\theta|<\pi/2
 \text{ on }X\}.
\]
When $\cH\ne\varnothing$, choose $\alpha$ within its class so that
$0\in\cH$. For $\varphi\in\cH$, put
\begin{align}
 \rho_\varphi&=\operatorname{Re}\left(e^{-\sqrt{-1}\widehat\theta}
               (\omega+\sqrt{-1}\alpha_\varphi)^n\right)>0,\nonumber\\
 \sigma_\varphi&=\operatorname{Im}\left(e^{-\sqrt{-1}\widehat\theta}
               (\omega+\sqrt{-1}\alpha_\varphi)^n\right).
 \label{phase-volume-and-error}
\end{align}
As in \cite{CY18}, define $\mathcal J\colon\cH\to\R$ by
\[
 d\mathcal J_\varphi(v)=-\int_Xv\sigma_\varphi,\qquad \mathcal J(0)=0.
\]
For the metric
\[
 \langle v,w\rangle_\varphi=\int_Xvw\rho_\varphi,
 \qquad v,w\in C^\infty(X,\R),
\]
the negative gradient flow is
\[
 \partial_t\varphi=-\operatorname{grad}\mathcal J(\varphi)
 =\frac{\sigma_\varphi}{\rho_\varphi}
 =\tan\bigl(\vartheta_\omega(\alpha_\varphi)-\widehat\theta\bigr).
\]
Along this flow,
\begin{equation}\label{tangent-flow-energy-identity}
 \frac{d}{dt}\mathcal J(\varphi(t))
 =-\int_X\tan^2\bigl(\vartheta_\omega(\alpha_\varphi)-\widehat\theta\bigr)
 \rho_\varphi.
\end{equation}

Let $d_1$ be the path-length distance induced by
$\|v\|_{1,\varphi}=\int_X|v|\rho_\varphi$. For $0<\theta<\pi/2$, set
\[
 \cH_\varepsilon=\{\varphi\in\cH:
 Q_\omega(\alpha_\varphi)>\varepsilon\text{ on }X\},
 \qquad 0<\varepsilon<\theta.
\]

\begin{theoremrecall}{thm-properness}
\PropernessStatement
\end{theoremrecall}

\begin{enumerate}
\item \emph{Tangent flow.}
\[
 \partial_t\varphi=
 \tan\bigl(\vartheta_\omega(\alpha_\varphi)-\widehat\theta\bigr).
\]
For $0<\theta<\pi/2$, long-time existence holds from every
$\varphi_0\in\cH$. A smooth $C$-subsolution gives convergence
\cite[Theorem 1.4]{Tak20}.

\item \emph{Cotangent flow.}
\[
 \partial_t\varphi=\cot Q_\omega(\alpha_\varphi)-\cot\theta.
\]
For $0<\theta<\pi$, long-time existence holds when
$Q_\omega(\alpha_{\varphi_0})<\pi$ on $X$. Convergence follows from a
smooth $C$-subsolution with $Q_\omega(\alpha_{\underline\varphi})<\pi$
\cite[Theorems 1.2--1.3]{FYZ25}.

\item \emph{The original flow.}
\[
 \partial_t\varphi=\vartheta_\omega(\alpha_\varphi)-\widehat\theta.
\]
Finite-time singularities can occur even in solvable classes
\cite{CJ25}. Local convergence near a solution is known without a
phase restriction \cite{HJ20}.
\end{enumerate}

For the original flow with $0<\theta<\pi/2$, a hypercritical initial
potential which is itself a $C$-subsolution yields smooth convergence
\cite[Theorem 1.3, citing Collins--Jacob--Yau, Remark 7.4]{Tak20}.
There is also a symmetry-restricted global existence result: on
$\operatorname{Bl}_p\bbP^n$, $n\geq3$, Calabi-symmetric initial and
background forms with
\[
 \vartheta_\omega(\alpha_{\varphi_0})>(n-2)\pi/2
\]
produce a flow existing for all time \cite[Theorem 1.2]{CJ25}.
This theorem does not assume smooth solvability and does not by itself
assert smooth convergence.

\Needspace{13\baselineskip}
\begin{problem}\label{problem-original-flow-basin}
Let $(X,\omega)$ be a compact K\"ahler $n$-fold with $n\geq2$, let
$\alpha$ be a smooth closed real $(1,1)$-form, and fix
$(n-1)\pi/2<\widehat\theta<n\pi/2$. Assume both of the following:
\begin{itemize}
\item $\vartheta_\omega(\alpha+\ddbar\varphi_*)=\widehat\theta$ for
some $\varphi_*\in C^\infty(X,\R)$.
\item The initial potential $\varphi_0\in C^\infty(X,\R)$ satisfies
$|\vartheta_\omega(\alpha+\ddbar\varphi_0)-\widehat\theta|<\pi/2$
on $X$.
\end{itemize}
Does the evolution
\[
 \partial_t\varphi=\vartheta_\omega(\alpha+\ddbar\varphi)-\widehat\theta,
 \qquad \varphi(0)=\varphi_0,
\]
exist for all time and converge smoothly modulo constants?
\end{problem}

This question isolates the initial-data issue raised after
\cite[Theorem 1.3]{Tak20}; it is not the Thomas--Yau conjecture itself.
The tangent-flow analogue follows from \cite[Theorem 1.4]{Tak20}, since
a smooth solution is a $C$-subsolution. Moreover,
\[
 \widehat\theta>\frac{(n-1)\pi}{2},\quad \varphi_0\in\cH
 \quad\Longrightarrow\quad
 \vartheta_\omega(\alpha_{\varphi_0})
 >\widehat\theta-\frac\pi2>\frac{(n-2)\pi}{2}.
\]
The hypotheses therefore lie inside the supercritical region, but do
not require the initial potential to be hypercritical or a
$C$-subsolution. The cited local, symmetric and singularity results do
not prove or disprove this precise general statement.

\section{Surface stability}\label{surface-stability}
Let $X$ be a smooth projective surface, let
$\Omega\in\NS(X)_\R$ be ample, and let $A=c_1(L)$ for a line bundle
$L$. Write
\[
 v=\langle\Omega^2,[X]\rangle,\quad
 a=\langle A^2,[X]\rangle,\quad
 s=\langle A\smile\Omega,[X]\rangle>0,\quad
 c=\frac{v-a}{2s},\quad A_c=A+c\Omega.
\]
For $B\in\NS(X)_\R$, define the twisted slope of a torsion-free sheaf
$F$ by
\[
 \mu_{B,\Omega}(F)
 =\frac{\langle(c_1(F)-\operatorname{rk}(F)B)\smile\Omega,[X]\rangle}
        {\operatorname{rk}(F)}.
\]
Let $\mathcal T_{B,\Omega}$ be generated under extensions by torsion
sheaves and torsion-free slope-semistable sheaves of positive twisted
slope. Let $\mathcal F_{B,\Omega}$ be generated by the torsion-free
slope-semistable sheaves of nonpositive twisted slope. The tilted heart is
\[
 \mathcal A_{B,\Omega}
 =\left\{E\in D^b(X):
 \begin{array}{l}
 H^{-1}(E)\in\mathcal F_{B,\Omega},\quad
 H^0(E)\in\mathcal T_{B,\Omega},\\
 H^j(E)=0\quad(j\ne-1,0)
 \end{array}\right\}.
\]
The divisorial stability condition $\sigma_{B,\Omega}$ uses this heart
and the central charge
\[
 Z_{B,\Omega}(E)
 =-\langle e^{-(B+\sqrt{-1}\Omega)}\operatorname{ch}(E),[X]\rangle
 \qquad\text{\cite[Section 2]{Mak26}.}
\]
For $B=0$, in particular,
\begin{equation}\label{categorical-central-charge}
 Z_{0,\Omega}(L)=\frac{v-a}{2}+\sqrt{-1}s
               =\frac12Z_X(A,\Omega).
\end{equation}
A nonzero $E\in\mathcal A_{B,\Omega}$ is $\sigma_{B,\Omega}$-stable if
\[
 \operatorname{Arg}Z_{B,\Omega}(G)
 <\operatorname{Arg}Z_{B,\Omega}(E)
 \quad\text{for every }0\ne G\subsetneq E
 \text{ in }\mathcal A_{B,\Omega},
 \qquad\operatorname{Arg}\in(0,\pi].
\]
At $B=0$, $s>0$ gives $L\in\mathcal A_{0,\Omega}$ and
$\operatorname{Arg}Z_{0,\Omega}(L)=\widehat\theta$.

For an integral curve $C\subset X$, set
\begin{equation}\label{curve-invariant}
 N_C=\langle\Omega,[C]\rangle(v-a)+2s\langle A,[C]\rangle
     =2s\langle A_c,[C]\rangle.
\end{equation}
By Theorem~\ref{thm-surface-criterion}, dHYM solvability is equivalent
to $N_C>0$ for every integral curve $C\subset X$.

\Needspace{13\baselineskip}
\begin{theoremrecall}{thm-stability-under-scaling}
\SurfaceStabilityStatement
\end{theoremrecall}

Theorem~\ref{thm-stability-under-scaling} is from the August 2026
preprint \cite{Mak26}.

Simultaneous scaling gives
\[
 (A,\Omega)\longmapsto(kA,k\Omega),\qquad
 \vartheta_{k\omega}(k\alpha_\varphi)=\vartheta_\omega(\alpha_\varphi).
\]
For a divisor subobject $L^{\otimes k}(-C)\subset L^{\otimes k}$ in
$\mathcal A_{0,k\Omega}$, expansion of the central charges gives
\begin{align*}
 &\operatorname{Im}\left(
 Z_{0,k\Omega}(L^{\otimes k})\,
 \overline{Z_{0,k\Omega}(L^{\otimes k}(-C))}\right)\\
 &=\frac{k^3}{2}\left(N_C-\frac{C^2s}{k}\right).
\end{align*}
Consequently,
\[
 \operatorname{Arg}Z_{0,k\Omega}(L^{\otimes k}(-C))
 <\operatorname{Arg}Z_{0,k\Omega}(L^{\otimes k})
 \quad\Longleftrightarrow\quad N_C-\frac{C^2s}{k}>0.
\]
In particular,
\[
 C^2<0,\quad N_C=0
 \quad\Longrightarrow\quad
 N_C-\frac{C^2s}{k}>0\qquad(k\geq1).
\]
Thus the strict subobject inequality can hold for every $k\geq1$
while the dHYM inequality $N_C>0$ fails.

For the blow-up in Example~\ref{ex-nef-boundary-class}, take
$L=\mathcal O_X(H)$. Then $N_C\geq0$ for every integral curve and
$N_E=0$. Thus $L^{\otimes k}$ is
$\sigma_{0,k\Omega}$-stable for every integer $k\geq1$, but
$A_c=H\notin\cK_X$ \cite[Example 2.11]{Fan26}.

This is a boundary example: the numerical class is nef rather than
K\"ahler. With $M=L^\vee$, both the ample-dual hypothesis and the strict
genericity hypothesis of Theorem~\ref{thm-toric-mirror} fail.
It disproves the implication from the stated scaled Bridgeland
stability to smooth dHYM solvability, not every possible stability
interpretation of the equation.

\Needspace{10\baselineskip}
\begin{theorem}[M\"akel\"a {\cite[Theorem 1.1(a)]{Mak26}}]
\label{thm-effective-divisor-detector}
Let $X$ be a smooth projective surface, let $B,\Omega\in\NS(X)_\R$
with $\Omega$ ample, and let $L$ be a line bundle with
$s_B=\langle\Omega\smile(c_1(L)-B),[X]\rangle>0$.
Then $L$ is not $\sigma_{B,\Omega}$-stable if and only if there is a
nonzero effective divisor $D$ satisfying
\begin{itemize}
\item $D^2<0$ and $\langle\Omega,[D]\rangle<s_B$;
\item $\operatorname{Arg}Z_{B,\Omega}(L(-D))
       \geq\operatorname{Arg}Z_{B,\Omega}(L)$.
\end{itemize}
\end{theorem}
The divisor need not be reduced or irreducible. The following stronger
question is distinguished explicitly in \cite[Remark 1.2]{Mak26}.

\begin{problem}
\label{problem-integral-curve-detector}
Let $X$ be a smooth projective surface and fix $B,\Omega\in\NS(X)_\R$
with $\Omega$ ample. Use the central charge
\[
 Z_{B,\Omega}(\mathcal E)
 =-\big\langle e^{-(B+\sqrt{-1}\Omega)}
 \operatorname{ch}(\mathcal E),[X]\big\rangle.
\]
Let $L$ be a line bundle satisfying
\begin{itemize}
\item $s_B=\langle\Omega\smile(c_1(L)-B),[X]\rangle>0$;
\item $L$ is not $\sigma_{B,\Omega}$-stable.
\end{itemize}
Must there be an integral curve $C$ such that
\[
 C^2<0,\qquad \langle\Omega,[C]\rangle<s_B,\qquad
 \operatorname{Arg}Z_{B,\Omega}(L(-C))
 \geq\operatorname{Arg}Z_{B,\Omega}(L)?
\]
Here $\operatorname{Arg}$ takes values in $(0,\pi]$.
\end{problem}

For the effective divisor $D$ supplied by
Theorem~\ref{thm-effective-divisor-detector}, write
$D=\sum_i m_iC_i$. Then
\[
 D^2=\sum_i m_i^2C_i^2+
       2\sum_{i<j}m_im_j(C_i\cdot C_j).
\]
Since $C_i\cdot C_j\geq0$ for $i\ne j$, the inequality $D^2<0$
gives a component $C_i$ with
\[
 C_i^2<0,\qquad
 0<\langle\Omega,[C_i]\rangle\leq\langle\Omega,[D]\rangle<s_B.
\]
A component satisfying the phase inequality would answer
Problem~\ref{problem-integral-curve-detector}; that inequality does not
follow from this intersection calculation. Thus the effective-divisor
criterion settles the unqualified divisorial formulation, whereas the
integral-curve question is an additional refinement, not a substitute
for the existence problem in Section~\ref{Thomas-Yau}.

\begin{thebibliography}{CJY20}
\bibitem[CJY20]{CJY20} T. C. Collins, A. Jacob and S.-T. Yau, \emph{$(1,1)$ forms with specified Lagrangian phase: A priori estimates and algebraic obstructions}, Camb. J. Math. \textbf{8} (2020), no. 2, 407--452. The theorem and proof numbering used here is that of \href{https://arxiv.org/abs/1508.01934v1}{arXiv:1508.01934v1} (2015).

\bibitem[CLT24]{CLT24} J. Chu, M.-C. Lee and R. Takahashi, \emph{A Nakai--Moishezon type criterion for supercritical deformed Hermitian--Yang--Mills equation}, J. Differential Geom. \textbf{126} (2024), no. 2, 583--632. Cited preprint version: \href{https://arxiv.org/abs/2105.10725v3}{arXiv:2105.10725v3} (2021).

\bibitem[JY17]{JY17} A. Jacob and S.-T. Yau, \emph{A special Lagrangian type equation for holomorphic line bundles}, Math. Ann. \textbf{369} (2017), 869--898. Cited preprint version: \href{https://arxiv.org/abs/1411.7457v1}{arXiv:1411.7457v1} (2014). \href{https://doi.org/10.1007/s00208-016-1467-1}{doi:10.1007/s00208-016-1467-1}.

\bibitem[Zha23]{Zha23} J. Zhang, \emph{A note on the supercritical deformed Hermitian--Yang--Mills equation}, \href{https://arxiv.org/abs/2302.06592v6}{arXiv:2302.06592v6} (2023), preprint.

\bibitem[DP04]{DP04} J.-P. Demailly and M. P\u{a}un, \emph{Numerical characterization of the K\"ahler cone of a compact K\"ahler manifold}, Ann. of Math. \textbf{159} (2004), 1247--1274. \href{https://annals.math.princeton.edu/2004/159-3/p05}{Publisher record and article}.

\bibitem[Yau78]{Yau78} S.-T. Yau, \emph{On the Ricci curvature of a compact K\"ahler manifold and the complex Monge--Amp\`ere equation, I}, Comm. Pure Appl. Math. \textbf{31} (1978), 339--411. \href{https://doi.org/10.1002/cpa.3160310304}{doi:10.1002/cpa.3160310304}.

\bibitem[Lin23]{Lin23} C.-M. Lin, \emph{On the Solvability of General Inverse $\sigma_k$ Equations}, \href{https://arxiv.org/abs/2310.05339v1}{arXiv:2310.05339v1} (2023), preprint; especially Theorem 1.5.

\bibitem[CY18]{CY18} T. C. Collins and S.-T. Yau, \emph{Moment maps, nonlinear PDE, and stability in mirror symmetry}, \href{https://arxiv.org/abs/1811.04824v2}{arXiv:1811.04824v2} (2018), extended preprint. Related published part: \emph{Moment maps, nonlinear PDE, and stability in mirror symmetry, I: Geodesics}, Ann. PDE \textbf{7} (2021), article 11. \href{https://doi.org/10.1007/s40818-021-00100-7}{doi:10.1007/s40818-021-00100-7}.

\bibitem[CL26]{CL26} J. Chu and M.-C. Lee, \emph{Hypercritical deformed Hermitian--Yang--Mills equation}, Trans. Amer. Math. Soc. \textbf{379} (2026), no. 8, 5765--5808. \href{https://doi.org/10.1090/tran/9655}{doi:10.1090/tran/9655}. The properness statement cited here is Theorem 1.4 of \href{https://arxiv.org/abs/2107.13192v1}{arXiv:2107.13192v1} (2021), with the theorem numbering and stated range of that version.

\bibitem[Tak20]{Tak20} R. Takahashi, \emph{Tan-concavity property for Lagrangian phase operators and applications to the tangent Lagrangian phase flow}, \href{https://arxiv.org/abs/2002.05132v6}{arXiv:2002.05132v6} (2020). Published in Int. J. Math.; \href{https://doi.org/10.1142/S0129167X20501165}{doi:10.1142/S0129167X20501165}.

\bibitem[FYZ25]{FYZ25} J. Fu, S.-T. Yau and D. Zhang, \emph{A new flow solving the LYZ equation in K\"ahler geometry}, \href{https://arxiv.org/abs/2105.13576v5}{arXiv:2105.13576v5} (2025; first version 2021). The arXiv record reports publication in J. Differential Geom.; theorem locators here refer to v5, whose running title is \emph{A new flow in K\"ahler geometry}.

\bibitem[CJ25]{CJ25} Y. H. Chan and A. Jacob, \emph{Singularity formation along the line bundle mean curvature flow}, Int. Math. Res. Not. \textbf{2025}, no. 5, rnaf037. \href{https://doi.org/10.1093/imrn/rnaf037}{doi:10.1093/imrn/rnaf037}. Cited preprint version: \href{https://arxiv.org/abs/2310.17709v1}{arXiv:2310.17709v1} (2023).

\bibitem[HJ20]{HJ20} X. Han and X. Jin, \emph{Stability of line bundle mean curvature flow}, \href{https://arxiv.org/abs/2001.07406v2}{arXiv:2001.07406v2} (2020); Theorem 1.5.

\bibitem[Mak26]{Mak26} A. M\"akel\"a, \emph{Negative Effective Divisors and Bridgeland Stability of Line Bundles on Surfaces}, \href{https://arxiv.org/abs/2608.26080v1}{arXiv:2608.26080v1} (26 August 2026), preprint.

\bibitem[Fan26]{Fan26} Y.-W. Fan, \emph{Notes on the deformed Hermitian--Yang--Mills equations and the large scaling limits of stability conditions}, \href{https://arxiv.org/abs/2604.22246v1}{arXiv:2604.22246v1} (24 April 2026), preprint.


\bibitem[TY02]{TY02} R. P. Thomas and S.-T. Yau,
\emph{Special Lagrangians, stable bundles and mean curvature flow},
Comm. Anal. Geom. \textbf{10} (2002), no. 5, 1075--1113.
Cited version: \href{https://arxiv.org/abs/math/0104197v3}{arXiv:math/0104197v3};
especially Sections 1--3, conditions (7.1)--(7.2), and Conjecture 7.3.

\bibitem[Nev13]{Nev13} A. Neves,
\emph{Finite time singularities for Lagrangian mean curvature flow},
Ann. of Math. \textbf{177} (2013), no. 3, 1029--1076.
Cited version: \href{https://arxiv.org/abs/1009.1083v3}{arXiv:1009.1083v3}
(2012), Theorem 6.1 and Remark 1.3.
\href{https://doi.org/10.4007/annals.2013.177.3.5}{doi:10.4007/annals.2013.177.3.5}.

\bibitem[LO24]{LO24} J. D. Lotay and G. Oliveira,
\emph{Special Lagrangians, Lagrangian mean curvature flow and the
Gibbons--Hawking ansatz},
J. Differential Geom. \textbf{126} (2024), no. 3, 1121--1184.
Cited version: \href{https://arxiv.org/abs/2002.10391v2}{arXiv:2002.10391v2}
(2022), Theorem 1.2 and Definition 6.2.
\href{https://doi.org/10.4310/jdg/1717348872}{doi:10.4310/jdg/1717348872}.

\bibitem[Joy15]{Joy15} D. Joyce,
\emph{Conjectures on Bridgeland stability for Fukaya categories of
Calabi--Yau manifolds, special Lagrangians, and Lagrangian mean curvature flow},
EMS Surv. Math. Sci. \textbf{2} (2015), 1--62.
Cited version: \href{https://arxiv.org/abs/1401.4949v2}{arXiv:1401.4949v2}
(2014), Conjectures 3.2 and 3.34.
\href{https://doi.org/10.4171/EMSS/8}{doi:10.4171/EMSS/8}.

\bibitem[Sto25]{Sto25} J. Stoppa,
\emph{Nakai--Moishezon criteria and the toric Thomas--Yau conjecture},
\href{https://arxiv.org/abs/2505.07228v2}{arXiv:2505.07228v2} (2025),
preprint. Theorem 1.2, Theorem 2.5, Corollary 2.7 and
Remark 2.6(v); the genericity condition is retained from this version.

\end{thebibliography}

\end{document}
